% Abstract, reader note and Figure 1.
\begin{ezsynnote}
\noindent{\sffamily\sbseries\color{ezsyn}How to read this draft.}\enspace This is the paper we intend
to earn, written before its decisive experiments. Framing, method, protocol and historical evidence
are real; \textbf{every EurekaZero result is synthetic}, set in \textcolor{ezsyn}{purple} and linked
to the gap ledger (\cref{app:gaps}), which names the experiment that must replace it. Figures with
synthetic data carry a \textsf{SYNTHETIC} tag; real results (\cref{app:history}) carry \est{} and
\pre{} badges.
\end{ezsynnote}

\begin{ezabstract}
\noindent{\sffamily\sbseries\small\color{ezaccent}Abstract}\par\smallskip
\noindent AlphaGo Zero showed that a network trained on the outcomes of its own search can play
strong moves without searching. We argue that innovation admits the same treatment. When a
problem's requirements can be encoded, its candidate solutions written in a typed language and
their quality measured by an independent verifier, inventing a solution is a move in a game against
the best solution known so far. We formalize this \emph{innovation game} and introduce EurekaZero,
which maps an encoded problem state through a learned solution embedding to one typed, executable
candidate (\code{game\_state\_encoding} $\to$ \code{solution\_embedding} $\to$ \code{solution}) and
is trained on verified outcomes of earlier searches. \emph{Zero} is a protocol: one candidate is
committed before any candidate-specific execution, scoring or revision. On sealed problem families
in three domains---learning procedures for unseen datasets, next-experiment decisions from
optimization histories, and load-bearing parts from engineering specifications---one EurekaZero
proposal improves failure-inclusive utility over the strongest predeclared non-search control by
\syn[EZ-G01]{\EZdA}, \syn[EZ-G03]{\EZdB} and \syn[EZ-G05]{\EZdC}, matches search after
\syn[EZ-G02]{\EZbeqA} and \syn[EZ-G06]{\EZbeqC} evaluated candidates, and beats frontier language
models given the same information and one attempt. Ablations attribute the gain to outcome training
and the solution embedding rather than to retrieval or catalog selection, and one shared model plays
all three games. Of committed learning procedures, \syn[EZ-G15]{\EZnovAfour\%} are useful and absent
from the frozen training archive; \syn[EZ-G15]{\EZnFive} proposed mechanisms survive removal and
replacement on reserved problems, \syn[EZ-G15]{\EZnSix} of them independently replicated. We release
the protocol, the benchmarks and every committed proposal.
\end{ezabstract}

\begin{figure}[t]
\centering
% Figure 1 — innovation as a game: Go and three innovation games, row by row. Conceptual; no data.
\begin{tikzpicture}[x=1cm, y=1cm,
  cell/.style={font=\sffamily\scriptsize, text width=3.05cm, align=center, inner sep=1pt, text=ezink},
  rowlab/.style={font=\sffamily\sbseries\scriptsize, color=ezaccent, anchor=west, text width=1.25cm, align=left},
  head/.style={font=\sffamily\sbseries\small, text=ezink, align=center, text width=3.1cm}]
  % column backdrops
  \fill[ezpaper, rounded corners=4pt] (1.42,1.05) rectangle (4.58,-6.32);
  \fill[ezaccentsoft, rounded corners=4pt] (4.72,1.05) rectangle (14.88,-6.32);
  \node[font=\sffamily\scriptsize, text=ezgrey, anchor=south] at (3.0,1.05) {precedent};
  \node[font=\sffamily\sbseries\scriptsize, text=ezaccent, anchor=south] at (9.8,1.05) {innovation games played by EurekaZero};

  % ---------------- icons
  % Go board
  \begin{scope}[shift={(3.0,0.1)}]
    \fill[ezamber!22, rounded corners=1pt] (-0.72,-0.72) rectangle (0.72,0.72);
    \foreach \i in {0,...,6} {
      \draw[ezgrey!80, line width=0.3pt] (-0.6+0.2*\i,-0.6) -- (-0.6+0.2*\i,0.6);
      \draw[ezgrey!80, line width=0.3pt] (-0.6,-0.6+0.2*\i) -- (0.6,-0.6+0.2*\i);}
    \foreach \p in {(-0.2,0.0),(0.2,0.2),(-0.4,0.4),(0.0,-0.4),(0.4,-0.2)} \fill[ezink] \p circle (0.082);
    \foreach \p in {(0.0,0.2),(-0.2,-0.2),(0.2,-0.4),(0.4,0.4)} \filldraw[fill=white, draw=ezink, line width=0.3pt] \p circle (0.082);
    \fill[ezink] (0.2,0.0) circle (0.082);
    \draw[ezwarm, line width=0.9pt] (0.2,0.0) circle (0.14);
  \end{scope}
  % Learning procedure: data table -> program
  \begin{scope}[shift={(6.4,0.1)}]
    \foreach \r in {0,...,5} \foreach \c in {0,...,4} {
      \pgfmathsetmacro\shade{\c==4 ? 1 : 0}
      \ifdim\shade pt=1pt \filldraw[fill=ezaccent!35, draw=white, line width=0.4pt] (-1.25+0.2*\c,-0.6+0.2*\r) rectangle ++(0.2,0.2);
      \else \filldraw[fill=white, draw=ezrule, line width=0.4pt] (-1.25+0.2*\c,-0.6+0.2*\r) rectangle ++(0.2,0.2); \fi}
    \draw[ezarrow] (-0.2,0.0) -- (0.18,0.0);
    \node[draw=ezaccent, fill=white, rounded corners=1.5pt, inner sep=2.5pt, font=\ttfamily\fontsize{5.4pt}{6.2pt}\selectfont, align=left, anchor=west]
      at (0.24,0.0) {fn fit(\ldots)\\fn predict(\ldots)\\\{ \ldots \}};
  \end{scope}
  % Next experiment: observations and the proposed point
  \begin{scope}[shift={(9.8,0.1)}]
    \draw[ezrule, line width=0.5pt, -{Stealth[length=1.4mm]}] (-0.95,-0.62) -- (0.98,-0.62);
    \draw[ezrule, line width=0.5pt, -{Stealth[length=1.4mm]}] (-0.95,-0.62) -- (-0.95,0.7);
    \draw[ezgrey!60, dashed, line width=0.5pt] plot[smooth, domain=-0.9:0.9, samples=40] (\x, {0.42*cos(deg(2.6*\x))-0.05*\x});
    \foreach \p in {(-0.78,0.0),(-0.45,0.17),(-0.1,-0.23),(0.36,-0.08),(0.72,0.16)} \fill[ezink] \p circle (0.05);
    \node[star, star points=5, star point ratio=2.2, fill=ezwarm, inner sep=0pt, minimum size=0.26cm] at (0.12,0.44) {};
    \node[font=\sffamily\fontsize{5.6pt}{6pt}\selectfont, text=ezwarm, anchor=west] at (0.22,0.47) {$x_{t+1}$};
  \end{scope}
  % Engineering part: bracket under load
  \begin{scope}[shift={(13.2,0.05)}]
    \filldraw[fill=ezgrey!25, draw=ezink, line width=0.5pt]
      (-0.85,-0.62) -- (0.85,-0.62) -- (0.85,-0.40) -- (-0.30,-0.40) -- (-0.58,0.12) -- (-0.58,0.62) -- (-0.85,0.62) -- cycle;
    \fill[white] (0.55,-0.51) circle (0.055); \fill[white] (0.15,-0.51) circle (0.055); \fill[white] (-0.715,0.40) circle (0.055);
    \draw[ezwarm, line width=0.9pt, -{Stealth[length=1.8mm]}] (-0.1,0.78) -- (-0.52,0.42);
    \node[font=\sffamily\fontsize{5.6pt}{6pt}\selectfont, text=ezwarm] at (0.05,0.86) {$F$};
    \fill[pattern=north east lines, pattern color=ezgrey] (-0.9,-0.70) rectangle (0.9,-0.64);
  \end{scope}

  % ---------------- column heads
  \node[head] at (3.0,-1.02) {Go\\[-1pt]{\scriptsize\mdseries\color{ezgrey}AlphaGo Zero}};
  \node[head] at (6.4,-1.02) {Learning procedure\\[-1pt]{\scriptsize\mdseries\color{ezgrey}domain A \textperiodcentered{} artifact}};
  \node[head] at (9.8,-1.02) {Next experiment\\[-1pt]{\scriptsize\mdseries\color{ezgrey}domain B \textperiodcentered{} action}};
  \node[head] at (13.2,-1.02) {Engineering part\\[-1pt]{\scriptsize\mdseries\color{ezgrey}domain C \textperiodcentered{} artifact}};

  % ---------------- rows
  \foreach \lab/\y in {State/-1.95, Move/-2.80, Rules/-3.65, Verifier/-4.50, Opponent/-5.35} {
    \node[rowlab] at (0.0,\y) {\lab};
    \draw[ezrule, line width=0.3pt] (1.5,\y+0.43) -- (14.8,\y+0.43);}
  \node[cell] at (3.0,-1.95) {a $19\times19$ position};
  \node[cell] at (6.4,-1.95) {training rows, objective, resource budget};
  \node[cell] at (9.8,-1.95) {evaluated points $(x_i,y_i)$, bounds, horizon};
  \node[cell] at (13.2,-1.95) {loads, interfaces, envelope, material, process};
  \node[cell] at (3.0,-2.80) {place one stone};
  \node[cell] at (6.4,-2.80) {commit one fit/predict program};
  \node[cell] at (9.8,-2.80) {commit one next query $x_{t+1}$};
  \node[cell] at (13.2,-2.80) {commit one CAD construction program};
  \node[cell] at (3.0,-3.65) {the rules of Go};
  \node[cell] at (6.4,-3.65) {typed procedure language, six guarantees};
  \node[cell] at (9.8,-3.65) {typed search space and budget};
  \node[cell] at (13.2,-3.65) {typed CAD language, validity checks};
  \node[cell] at (3.0,-4.50) {the game result};
  \node[cell] at (6.4,-4.50) {fit, then error on held-out rows};
  \node[cell] at (9.8,-4.50) {the objective's measured value};
  \node[cell] at (13.2,-4.50) {build, mesh, finite elements, load test};
  \node[cell] at (3.0,-5.35) {the previous network};
  \node[cell] at (6.4,-5.35) {incumbent procedure for that dataset};
  \node[cell] at (9.8,-5.35) {incumbent decision for that history};
  \node[cell] at (13.2,-5.35) {incumbent design for that specification};

  % ---------------- the game, in one line
  \node[font=\sffamily\scriptsize, text=ezaccent, align=center] at (9.8,-6.05)
    {state $\;\to\;$ one committed move $\;\to\;$ verified outcome $\;\to\;$ win if it beats the incumbent};
\end{tikzpicture}

\caption{Innovation as a game. The game AlphaGo Zero mastered (left) and three innovation games
EurekaZero plays share one structure. A \emph{state} encodes everything the player may see; a
\emph{move} is one complete, typed candidate committed before it is tested; the \emph{rules} are the
typed solution language and its guarantees; an independent \emph{verifier} scores the move after
commitment; and the \emph{opponent} is the incumbent, the best verified solution known for that
problem. Winning a game raises the bar for the next one. Unlike Go, each innovation game is a new
board drawn from a family the player has never seen.}
\label{fig:game}
\end{figure}
